\begin{figure}[htbp]
\centering
\setlength{\unitlength}{1pt}
\begin{picture}(390,150)
\put(0,136){\makebox(390,12)[l]{\small One historical connecting chain}}
\put(55,106){\line(1,0){230}}
\put(55,106){\circle*{5}}
\put(285,106){\circle*{5}}
\put(170,106){\makebox(0,0){$\times$}}
\put(55,86){\makebox(0,12){$x_-$}}
\put(285,86){\makebox(0,12){$x_+$}}
\put(170,117){\makebox(0,12){\small vanished connector}}
\put(0,66){\makebox(390,12)[l]{\small Preliminary survivors and fixed hull}}
\linethickness{1.4pt}
\put(55,42){\line(1,0){230}}
\linethickness{.4pt}
\put(55,42){\circle*{5}}
\put(285,42){\circle*{5}}
\put(145,27){\line(0,1){30}}
\put(345,27){\line(0,1){30}}
\put(145,7){\makebox(0,12){\small skipped line}}
\put(345,7){\makebox(0,12){\small not blocked here}}
\end{picture}
\caption{An inherited class viewed in a cut coordinate. A connector that
vanished before the preliminary tests no longer appears among the
surviving labels, but its connecting segments remain covered by the
stored train history. The surviving extremes therefore determine a
hull in which sampled raw lines are skipped. The outside candidate is
not blocked by this hull; other hulls are not shown. The preliminary
hull is fixed before the lines are sampled and is not recomputed after
later purges. Only one possible chain and its class are shown;
historical chains may backtrack and different classes' hulls may overlap.}
\label{fig:inherited-hull}
\end{figure}
